\documentclass[final]{beamer}
\usepackage[size=custom,width=90,height=142,scale=1.25]{beamerposter}
\usetheme{default}
\usepackage{fontspec}
\setmainfont{Nimbus Sans} % Sustituta compatible con Helvetica disponible en el sistema.
\setsansfont{Nimbus Sans}
\usepackage[spanish,es-nodecimaldot]{babel}
\usepackage{graphicx,booktabs,array,tikz,xcolor}
\usepackage[most]{tcolorbox}
\usepackage{ragged2e}
\usepackage{hyperref}
% En el repositorio: poster/ junto a resultados/. El segundo path permite
% compilar esta copia de preparación desde el directorio de trabajo actual.
\graphicspath{{../resultados/}{../../../3. Metodologia/BPN/resultados/}}
\definecolor{navy}{HTML}{4B658E}
\definecolor{blue}{HTML}{3F6291}
\definecolor{teal}{HTML}{2A9D9C}
\definecolor{violet}{HTML}{7774D6}
\definecolor{pale}{HTML}{DCE8E6}
\definecolor{ink}{HTML}{2B303A}
\definecolor{muted}{HTML}{5A6270}
\definecolor{headerblue}{HTML}{4B658E}
\definecolor{headerteal}{HTML}{2A9D9C}
\definecolor{sidebarbg}{HTML}{DCE8E6}
\definecolor{titleblue}{HTML}{3A5A8C}
\definecolor{textdark}{HTML}{2B303A}
\definecolor{textmuted}{HTML}{5A6270}
\definecolor{barbg}{HTML}{CBD5E1}
\definecolor{barfill}{HTML}{3F6291}
\setbeamercolor{normal text}{fg=ink,bg=white}
\setbeamercolor{block title}{fg=white,bg=navy}
\setbeamertemplate{navigation symbols}{}
\setbeamertemplate{itemize items}[circle]
\setlength{\columnsep}{1.1cm}
\setlength{\parskip}{0.22em}
\newcommand{\sectionbar}[2]{\par\noindent{\color{titleblue}\large\bfseries #1\hspace{0.25em}#2}\par\vspace{0.10cm}{\color{barfill}\hrule height 0.65pt}\par\vspace{0.50cm}}
\newtcolorbox{sidebox}[1][]{colback=pale,colframe=pale,boxrule=0pt,arc=2.5mm,left=10mm,right=10mm,top=11mm,bottom=11mm,before skip=0pt,after skip=0pt,fontupper=\normalsize,width=0.89\linewidth,center,#1}
\newtcolorbox{resultbox}[1][]{colback=white,colframe=sidebarbg,boxrule=0.8pt,arc=0.7mm,left=3mm,right=3mm,top=2mm,bottom=2mm,fontupper=\normalsize,#1}
\newcommand{\chart}[2]{\begin{center}\includegraphics[width=\linewidth,height=#1,keepaspectratio]{#2}\end{center}}
\newcommand{\smallchart}[2]{\begin{center}\includegraphics[width=\linewidth,height=#1,keepaspectratio]{#2}\end{center}}
\begin{document}
\begin{frame}[t]
\begin{tcolorbox}[enhanced,colback=navy,colframe=navy,interior style={left color=headerblue,right color=headerteal},boxrule=0mm,arc=1mm,left=5mm,right=5mm,top=4mm,bottom=4mm]
\begin{columns}[c,totalwidth=\linewidth]
\begin{column}{0.28\linewidth}
\begin{minipage}[c]{0.25\linewidth}
\centering\includegraphics[width=4.6cm]{logo-udc.png}
\end{minipage}\hfill
\begin{minipage}[c]{0.71\linewidth}
\color{white}\raggedright
{\fontsize{39}{23}\selectfont\bfseries Universidad de Cartagena\par}
\vspace{1mm}
{\fontsize{29}{20}\selectfont Programa de economía\par}
{\fontsize{29}{20}\selectfont Economía de la salud 2026-2\par}
\end{minipage}
\end{column}
\begin{column}{0.70\linewidth}
\centering\color{white}
\makebox[\linewidth][c]{\hspace*{-1cm}{\fontsize{70}{62}\selectfont\bfseries Bajo peso al nacer en Colombia}\hspace*{1cm}}\\[0.2cm]
{\fontsize{21}{25}\selectfont Tendencias, edad gestacional y desigualdades territoriales, 2005–2024\hspace{0.6em}·\hspace{0.6em}Estudio poblacional con Estadísticas Vitales del DANE\hspace{0.6em}·\hspace{0.6em}Colombia}\\[0.3cm]
{\fontsize{25}{20}\selectfont
  \textbf{Autores:} Luis Armando Ramírez Bárcenas · Edinson Polo Arrieta \hspace{0.45cm}
\textbf{Tutor:} Nelson Alvis Guzmán \hspace{0.45cm}
}\\[0.15cm]
{\fontsize{25}{20}\selectfont
\textbf{Colaboradores:} Michel Álvarez Paternina · Jean Hernández Herazo
}
\end{column}
\end{columns}
\end{tcolorbox}
\vspace{0.35cm}
\begin{columns}[t,totalwidth=\linewidth]
% Lateral: resumen, método y lectura
\begin{column}{0.38\linewidth}
\begin{sidebox}
\sectionbar{}{INTRODUCCIÓN}
El bajo peso al nacer (BPN, menos de 2.500 g) se asocia con consecuencias duraderas en salud, educación e ingresos. En Colombia, el BPN total aumentó mientras el BPN de nacidos a término disminuyó.
\vspace{0.12cm}
\sectionbar{}{OBJETIVO}
Analizar la evolución de las desigualdades socioeconómicas y territoriales del BPN en Colombia entre 2005 y 2024.
\vspace{0.12cm}
\sectionbar{}{DATOS Y MÉTODOS}
\begin{itemize}\setlength{\itemsep}{0.28em}
\item 12,9 millones de nacidos vivos del DANE, 2005–2024; 33 unidades: 32 departamentos y Bogotá D.C.
\item Tendencia anual y descomposición de Kitagawa por semana de gestación.
\item Razones de prevalencia departamentales frente a cinco departamentos de referencia; índice de Theil y brechas, crudos y ajustados por altitud y edad materna.
\item Regresión de Poisson e índice relativo de desigualdad (RII) para educación y afiliación materna.
\item Índice de concentración ponderado por nacidos vivos, con bootstrap por departamento.
\end{itemize}
\end{sidebox}
\vspace{1.1cm}
\begin{sidebox}
\sectionbar{}{HALLAZGOS CLAVE}
\centering{\color{teal}\fontsize{39}{43}\selectfont\bfseries +2,97}\par
\large puntos porcentuales en la prevalencia nacional\\
\normalsize 8,29\% (2005) → 11,25\% (2024)
\vspace{0.1cm}
\begin{itemize}\setlength{\itemsep}{0.28em}
\item El BPN aumentó en 30 de los 33 departamentos.
\item La proporción de nacimientos en semana 37 pasó de 9,5\% a 16,7\% (2008–2024).
\item Frente a educación superior, la RP fue 1,17 para BPN y 1,60 para BPN a término en madres sin educación o preescolar.
\item La brecha educativa relativa creció cerca de 10\% por década.
\end{itemize}
\end{sidebox}
\vspace{1.1cm}
\begin{sidebox}
\sectionbar{}{INTERPRETACIÓN}
Las diferencias territoriales se reducen al ajustar por altitud y edad materna. La concentración observada favorece a departamentos de mejor posición; tras el ajuste, el índice baja y es compatible con cero.
\end{sidebox}
\end{column}

% Área principal: resultados gráficos grandes, en el orden de la historia
\begin{column}{0.60\linewidth}
\sectionbar{1}{EVOLUCIÓN NACIONAL}
\begin{resultbox}
\centering La prevalencia subió de 8,29\% a 11,25\%; cuatro quintas partes del aumento ocurrieron desde 2015. La tendencia se mantiene al estandarizar por edad materna.
\end{resultbox}
\chart{32cm}{G1_bpn_nacional.png}
\vspace{-0.2cm}
\sectionbar{2}{EDAD GESTACIONAL: CRECE LA PROPORCIÓN EN SEMANA 37}
\chart{32cm}{G15_semanas_gestacion.png}
\vspace{-0.1cm}
\begin{resultbox}
Entre 2008 y 2024, la proporción de nacimientos en semana 37 pasó de 9,5\% a 16,7\%. Esta semana aportó 1,72 de los 2,45 puntos que subió el BPN; dentro de ella, el BPN pasó de 13,37\% a 17,96\%.
\end{resultbox}
\end{column}
\end{columns}
\vspace{0.2cm}
% Resultados 3 y 4 ocupan todo el ancho disponible, incluido el espacio bajo el lateral.
\begin{columns}[t,totalwidth=\linewidth]
\begin{column}{0.49\linewidth}
\sectionbar{3}{BRECHAS ENTRE DEPARTAMENTOS}
\smallchart{25cm}{G7_RP_departamentos_quinquenios.png}
En 2020–2024, Bogotá tuvo RP 2,15 frente a la referencia y 1,47 tras ajustar por altitud y edad. El ajuste reduce la dispersión a cerca de un tercio; La Guajira pasa al primer lugar. La desigualdad departamental no aumentó en veinte años.
\end{column}
\begin{column}{0.49\linewidth}
\sectionbar{4}{CONCENTRACIÓN SOCIOECONÓMICA}
\smallchart{25cm}{G11_curvas_estandarizadas_trienios.png}
El índice observado bajó de 0,089 (2008–2010) a 0,081 (2022–2024). Tras ajustar por altitud y edad materna, cae a 0,021 y es compatible con cero.
\end{column}
\end{columns}
\vspace{0.25cm}
\sectionbar{5}{EL GRADIENTE EDUCATIVO EXISTE Y SE AMPLÍA}
\begin{resultbox}[arc=2.5mm,left=6mm,right=6mm,top=3mm,bottom=3mm]
\begin{columns}[t,totalwidth=\linewidth]
\begin{column}{0.34\linewidth}
\begin{minipage}[t]{\linewidth}\vspace{0pt}
\RaggedRight\hyphenpenalty=10000\exhyphenpenalty=10000 Al comparar madres del mismo departamento y a igual altitud, una menor educación materna se asocia con más BPN, especialmente en nacidos a término.
\end{minipage}
\end{column}
\begin{column}{0.63\linewidth}
\begin{minipage}[t]{\linewidth}\vspace{0pt}
\centering
\begin{tabular*}{0.98\linewidth}{@{\extracolsep{\fill}}lcc@{}}
\toprule
Educación materna & BPN total & BPN a término \\
\midrule
Ninguna o preescolar & 1,17 (1,11–1,24) & 1,60 (1,44–1,77) \\
Primaria & 1,04 (1,00–1,09) & 1,30 (1,24–1,35) \\
Secundaria o media & 1,02 (0,99–1,05) & 1,15 (1,13–1,17) \\
\bottomrule
\end{tabular*}
\par\vspace{2mm}
\RaggedRight\small RP frente a educación superior (IC 95\%), 2008–2024. El gradiente relativo aumentó cerca de 10\% por década.
\end{minipage}
\end{column}
\end{columns}
\end{resultbox}
\vspace{-0.1cm}
\begin{tcolorbox}[colback=headerblue,colframe=headerblue,boxrule=0pt,arc=2.5mm,left=11mm,right=11mm,top=3mm,bottom=2mm]
\color{white}\makebox[\linewidth][c]{{\large\bfseries CONCLUSIONES}}\par\vspace{0.1cm}
\begin{columns}[t,totalwidth=\linewidth]
\begin{column}{0.23\linewidth}
El BPN aumentó casi 3 puntos; la semana 37 aportó 1,72 puntos del aumento desde 2008.
\end{column}
\begin{column}{0.23\linewidth}
Las diferencias departamentales dependen en parte de altitud y composición, y no crecieron en veinte años.
\end{column}
\begin{column}{0.23\linewidth}
Las madres con menor educación presentan más BPN, sobre todo a término, y la brecha se amplía.
\end{column}
\begin{column}{0.23\linewidth}
La desigualdad socioeconómica aparece al comparar madres; tras el ajuste, la concentración territorial es compatible con cero.
\end{column}
\end{columns}
\end{tcolorbox}
\vspace{-0.3cm}
\begin{columns}[t,totalwidth=\linewidth]
\begin{column}{0.49\linewidth}
\begin{tcolorbox}[colback=sidebarbg,colframe=sidebarbg,boxrule=0pt,arc=2.5mm,left=10mm,right=10mm,top=1.5mm,bottom=1.5mm]
{\footnotesize
\textbf{LIMITACIONES.} Estudio observacional: describe asociaciones, no causas. El peso se analiza agrupado; la semana 37 se separa solo desde 2008; el análisis territorial compara 33 unidades.\\[1mm]
\textbf{FUENTE.} DANE, Estadísticas Vitales, nacimientos 2005–2024. Cálculos propios. Código: \href{https://github.com/IntegralIndefinida/BPN}{github.com/IntegralIndefinida/BPN}.
}
\end{tcolorbox}
\end{column}
\begin{column}{0.49\linewidth}
\begin{tcolorbox}[colback=sidebarbg,colframe=sidebarbg,boxrule=0pt,arc=2.5mm,left=10mm,right=10mm,top=1.5mm,bottom=1.5mm]
{\footnotesize
\textbf{REFERENCIAS.}\\
Guarnizo-Herreño et al. (2021). PLOS ONE, 16(7), e0255150.\\
Silveira et al. (2019). International Journal of Epidemiology, 48(Supl. 1), i46–i53.\\
Kakwani et al. (1997). Journal of Econometrics, 77(1), 87–103.\\
Jensen y Moore (1997). American Journal of Public Health, 87(6), 1003–1007.
}
\end{tcolorbox}
\end{column}
\end{columns}
\end{frame}
\end{document}
